\begin{problem}[Exactness on a finite Sierpinski space]\label{prob:vi16-p18}
Describe sheaves of abelian groups on the Sierpinski space by a restriction map $A\to B$.  Translate exactness of a sequence of such sheaves into exactness of the two associated stalk sequences.
\end{problem}

\ifdefined\FullProblemDossiers
\paragraph{Why this problem matters.}
Finite spaces make the abstract stalkwise theorem visible as a finite diagram calculation.

\paragraph{Useful ideas.}
Compute the stalk at each of the two points from the restriction diagram.

\paragraph{Guided hints.}
\begin{enumerate}
\item Identify which open is the smallest neighborhood of each point.
\item The stalk at a point is the group attached to its smallest neighborhood.
\end{enumerate}

\begin{solution}
Let the opens be $\varnothing,\{x\},X$ with $X=\{\eta,x\}$.  A sheaf is encoded by $A=\F(X)$, $B=\F(\{x\})$, and a restriction $A\to B$.  The stalk at $\eta$ is $A$, while the stalk at $x$ is $B$.  Therefore a sequence of sheaves is exact exactly when the sequence of the $A$-groups and the sequence of the $B$-groups are both exact.
\end{solution}

\paragraph{What happened mathematically.}
The general theorem reduces to two ordinary exactness checks on this finite space.

\paragraph{Extensions.}
Repeat for a three-point specialization chain.

\paragraph{Related problems.}
VI/12 and VI/13 introduced finite-space presheaves and stalks.

\paragraph{Provenance.}
New synthesis.
\else
\paragraph{Provenance.} New synthesis.
\fi
